\documentclass[11pt,a4paper]{article}
\usepackage[utf8]{inputenc}
\usepackage[T1]{fontenc}
\usepackage{amsmath,amssymb,amsfonts}
\usepackage{booktabs}
\usepackage{longtable}
\usepackage{geometry}
\usepackage{graphicx}
\usepackage{xcolor}
\usepackage{enumitem}
\usepackage{titlesec}
\usepackage{fancyhdr}
\usepackage{hyperref}

\geometry{margin=2.2cm}
\hypersetup{colorlinks=true,linkcolor=blue!60!black,urlcolor=blue!60!black}

\pagestyle{fancy}
\fancyhf{}
\fancyhead[L]{\small RelaLeap: EPC Experiments R8 \& R9 Stage A}
\fancyhead[R]{\small \thepage}
\renewcommand{\headrulewidth}{0.4pt}

\titleformat{\section}{\large\bfseries}{\thesection}{1em}{}
\titleformat{\subsection}{\normalsize\bfseries}{\thesubsection}{1em}{}

\title{\textbf{Energy--Parameterized Credit-assignment (ePC):\\[4pt]
Experimental Report -- R8 Broad Outcome \& R9 Stage A\\[6pt]
Representation Battery}}
\author{ZeroBot Research Agent \\ \textit{for Benjamin Goertzel}}
\date{2026-07-22}

\begin{document}
\maketitle
\thispagestyle{fancy}

\begin{abstract}
This report documents two experiment batches on Energy--Parameterized
Credit-assignment (ePC), a training method that enforces an
energy--monotonicity property during knowledge distillation.
Experiment R8 evaluates raw predictive performance across 5 seeds and
5 arms at both matched-update and matched-wall-clock budgets.
Experiment R9 Stage A probes the learned representations along four
axes: calibration validity (A0), representational collapse (A1),
factor probe accuracy (A2), and causal selectivity (A3).
The headline finding is that ePC's energy--monotonicity property is
robustly validated (10/10 EPC records), but ePC does not outperform
plain baselines on validation loss, and its learned representations
exhibit significantly more collapse and faster loss of task-relevant
factor information.
\end{abstract}

\tableofcontents
\newpage

%% ================================================================
\section{Background and Motivation}
%% ================================================================

ePC (Energy--Parameterized Credit-assignment) is a training objective
that introduces an energy-based constraint into the knowledge
distillation (KD) process.  The core idea is that credit assignment
to deeper network layers should respect an energy--monotonicity
property: the ``activity energy'' of each layer should be
non-increasing along the inference path, providing a structured
alternative to ordinary backpropagated KD gradients.

Two variants are tested:
\begin{itemize}[nosep]
  \item \textbf{epc\_original}: standard ePC with default inference depth.
  \item \textbf{epc\_deep}: ePC with increased inference depth ($T{=}12$).
\end{itemize}

These are compared against three baselines:
\begin{itemize}[nosep]
  \item \textbf{bp\_ce}: ordinary backpropagation with cross-entropy loss.
  \item \textbf{bp\_kd}: ordinary backpropagation with KD loss, matched
    at 200 updates.
  \item \textbf{bp\_kd\_wallclock}: ordinary KD allowed to run for the
    same wall-clock budget as epc\_original ($\sim$445\,s), completing
    $\sim$6.7$\times$ more update steps.
\end{itemize}

The teacher model is a pretrained GPT-2 small (val\_loss 4.124,
val\_perplexity 61.8).  The student is a 6-layer GPT-2 trained on a
synthetic subject--verb--object grammar corpus.

%% ================================================================
\section{Experiment R8: Broad Outcome}
%% ================================================================

\subsection{Setup}

\begin{itemize}[nosep]
  \item \textbf{Pod:} RunPod A100 PCIe 80GB, \$1.39/hr
  \item \textbf{Source commit:} \texttt{ecf2f79}
  \item \textbf{Seeds:} 1729, 3253, 6421, 8191, 10103 (5 seeds)
  \item \textbf{Arms:} bp\_ce, bp\_kd, epc\_original, epc\_deep,
    bp\_kd\_wallclock (5 arms)
  \item \textbf{Total records:} 25 (5 seeds $\times$ 5 arms)
  \item \textbf{Matched-update arms:} 200 updates each (bp\_ce, bp\_kd,
    epc\_original, epc\_deep)
  \item \textbf{Matched-wall-clock arm:} bp\_kd\_wallclock runs for
    $\sim$445\,s (same elapsed budget as epc\_original), completing
    $\sim$1334 updates on average.
\end{itemize}

\subsection{R8 Results: Per-Seed Records}

\begin{longtable}{llrrrrr}
\toprule
\textbf{Arm} & \textbf{Seed} & \textbf{val\_loss} & \textbf{val\_ppl} & \textbf{kd\_gap} & \textbf{updates} & \textbf{elapsed (s)} \\
\midrule
\endhead
bp\_ce & 1729 & 6.265 & 525.7 & 1158.8 & 200 & 130 \\
bp\_ce & 3253 & 6.194 & 489.8 & 1164.6 & 200 & 64 \\
bp\_ce & 6421 & 6.222 & 503.9 & 1148.5 & 200 & 65 \\
bp\_ce & 8191 & 6.270 & 528.3 & 1163.6 & 200 & 65 \\
bp\_ce & 10103 & 6.353 & 574.0 & 1169.0 & 200 & 64 \\
\midrule
bp\_kd & 1729 & 7.053 & 1156.5 & 828.1 & 200 & 104 \\
bp\_kd & 3253 & 7.053 & 1155.8 & 830.3 & 200 & 66 \\
bp\_kd & 6421 & 6.950 & 1043.3 & 807.6 & 200 & 67 \\
bp\_kd & 8191 & 7.037 & 1137.5 & 818.6 & 200 & 67 \\
bp\_kd & 10103 & 7.103 & 1215.2 & 834.3 & 200 & 66 \\
\midrule
epc\_original & 1729 & 6.747 & 851.4 & 981.0 & 200 & 591 \\
epc\_original & 3253 & 6.656 & 777.7 & 987.9 & 200 & 406 \\
epc\_original & 6421 & 6.695 & 808.1 & 962.4 & 200 & 414 \\
epc\_original & 8191 & 6.777 & 877.5 & 992.0 & 200 & 409 \\
epc\_original & 10103 & 7.028 & 1128.3 & 966.5 & 200 & 403 \\
\midrule
epc\_deep & 1729 & 6.659 & 780.1 & 941.3 & 200 & 2032 \\
epc\_deep & 3253 & 6.615 & 746.5 & 953.6 & 200 & 1673 \\
epc\_deep & 6421 & 6.713 & 823.0 & 943.1 & 200 & 1670 \\
epc\_deep & 8191 & 6.869 & 961.8 & 970.9 & 200 & 1660 \\
epc\_deep & 10103 & 6.913 & 1004.9 & 994.1 & 200 & 1655 \\
\midrule
bp\_kd\_wallclock & 1729 & 5.587 & 267.0 & 418.9 & 1783 & 591 \\
bp\_kd\_wallclock & 3253 & 5.824 & 338.4 & 482.2 & 1220 & 406 \\
bp\_kd\_wallclock & 6421 & 5.711 & 302.3 & 465.6 & 1229 & 414 \\
bp\_kd\_wallclock & 8191 & 5.881 & 358.3 & 491.1 & 1225 & 409 \\
bp\_kd\_wallclock & 10103 & 5.918 & 371.7 & 491.0 & 1213 & 403 \\
\bottomrule
\caption{R8 per-seed per-arm results. Teacher: val\_loss=4.124, val\_ppl=61.8.}
\label{tab:r8-per-seed}
\end{longtable}

\subsection{R8 Results: Aggregate Statistics}

\begin{table}[h]
\centering
\begin{tabular}{lrrrrr}
\toprule
\textbf{Arm} & \textbf{val\_loss} (mean$\pm$sd) & \textbf{val\_ppl} & \textbf{kd\_gap} & \textbf{updates} & \textbf{elapsed (s)} \\
\midrule
bp\_ce & 6.261$\pm$0.060 & 524.3 & 1160.9 & 200 & 78 \\
bp\_kd & 7.039$\pm$0.056 & 1141.6 & 823.8 & 200 & 74 \\
epc\_original & 6.781$\pm$0.146 & 888.6 & 978.0 & 200 & 445 \\
epc\_deep & 6.754$\pm$0.131 & 863.2 & 960.6 & 200 & 1738 \\
bp\_kd\_wallclock & 5.785$\pm$0.135 & 327.5 & 469.8 & 1334 & 445 \\
\midrule
\textit{Teacher} & \textit{4.124} & \textit{61.8} & --- & --- & --- \\
\bottomrule
\end{tabular}
\caption{R8 aggregate statistics (5 seeds per arm).}
\label{tab:r8-aggregate}
\end{table}

\subsection{R8 Key Findings}

\begin{enumerate}
  \item \textbf{Energy--monotonicity holds robustly.}
    All 10 EPC records (5 seeds $\times$ \{epc\_original,
    epc\_deep\}) have \texttt{activity\_energy\_monotone: true}.
    This is the primary positive structural finding.

  \item \textbf{At matched update count (200), bp\_ce beats both EPC
    arms.}
    bp\_ce val\_loss 6.261 vs epc\_original 6.781 and epc\_deep 6.754.
    EPC does not improve over plain cross-entropy at equal updates.

  \item \textbf{At matched wall-clock, ordinary KD dominates.}
    bp\_kd\_wallclock completes $\sim$6.7$\times$ more updates (1334
    vs 200) and achieves val\_loss 5.785 / val\_ppl 327.5, far better
    than all 200-update arms.

  \item \textbf{epc\_deep is marginally better than epc\_original}
    on val\_loss (6.754 vs 6.781), but the difference is within seed
    variation ($\pm$0.13).  Increasing inference depth from $T{=}6$ to
    $T{=}12$ does not yield a clear performance gain.

  \item \textbf{bp\_kd at 200 updates is the worst arm}
    (val\_loss 7.039).  Plain KD at low update counts produces worse
    students than plain CE, consistent with the KD signal being noisy
    when undertrained.

  \item \textbf{EPC is substantially slower.}
    epc\_original averages 445\,s per run; epc\_deep averages 1738\,s
    (3.9$\times$ slower than epc\_original, 22$\times$ slower than
    bp\_ce).  The energy--monotonicity enforcement requires iterative
    inference at each step.
\end{enumerate}

\subsection{Credit Assignment Profile (EPC arms)}

The credit assignment diagnostic records the ePC gradient norm and
ordinary KD gradient norm per transformer block.  In both EPC arms,
the ePC gradient norm is effectively zero for blocks 0--4 and only
non-zero for the final block (h.5):

\begin{table}[h]
\centering
\begin{tabular}{llrrr}
\toprule
\textbf{Block} & \textbf{ePC norm} & \textbf{ordinary KD norm} & \textbf{norm ratio} & \textbf{finite} \\
\midrule
h.0 & 0.0000 & 28715.22 & 0.000000 & True \\
h.1 & 0.0000 & 22168.94 & 0.000000 & True \\
h.2 & 0.0000 & 18854.66 & 0.000000 & True \\
h.3 & 0.0000 & 16373.57 & 0.000000 & True \\
h.4 & 0.0000 & 14862.22 & 0.000000 & True \\
h.5 & 0.4913 & 15015.60 & 0.000033 & True \\
\bottomrule
\end{tabular}
\caption{Credit assignment for seed 3253, epc\_original.
  epc\_deep shows an identical pattern (ePC norm 0.4916 at h.5 only).}
\label{tab:credit-assignment}
\end{table}

This means the ePC signal is concentrated entirely in the last
transformer block, with the ratio ePC/KD $\approx 3.3 \times
10^{-5}$.  The energy--monotonicity constraint is structurally
enforced, but the credit it assigns is negligible relative to ordinary
KD gradients.

%% ================================================================
\section{Experiment R9 Stage A: Representation Battery}
%% ================================================================

\subsection{Setup}

R9 Stage A was run on the same pod immediately after R8 training
completed.  It applies four probe families to the 20 checkpoints from
the 4 matched-update arms (bp\_ce, bp\_kd, epc\_original, epc\_deep
$\times$ 5 seeds):

\begin{itemize}[nosep]
  \item \textbf{A0 -- Calibration:} CKA symmetry, self-CKA, known-rank
    recovery, permutation null accuracy.
  \item \textbf{A1 -- Collapse:} entropy effective rank, participation
    ratio, top-1 variance fraction, dead neuron fraction.  Measured at
    embedding, 6 transformer blocks, and final layer.
  \item \textbf{A2 -- Factor Probes:} linear classifier accuracy
    (test and held-out) for four syntactic factors: subj\_num,
    obj\_num, tense, negation.
  \item \textbf{A3 -- Causal Selectivity:} intervention-based
    selectivity score for the same four factors.
\end{itemize}

\subsection{A0: Calibration Validation}

\begin{table}[h]
\centering
\begin{tabular}{lr}
\toprule
\textbf{Metric} & \textbf{Value} \\
\midrule
CKA symmetry error & $6.55 \times 10^{-15}$ \\
CKA (X,Y) & 0.958 \\
CKA (self) & 1.000 \\
Known rank (input) & 5 \\
Known rank (recovered) & 5 \\
Permutation null train acc & 0.531 \\
Permutation null test acc & 0.448 \\
\bottomrule
\end{tabular}
\caption{A0 calibration (seed 3253, epc\_original).  All probes pass.}
\label{tab:a0-calibration}
\end{table}

CKA symmetry error is at machine-precision zero, self-CKA is exactly
1.0, and known-rank recovery is exact.  The probe infrastructure is
trustworthy.

%% ---- A1 ----
\subsection{A1: Representational Collapse}

This is the most diagnostically informative probe.  The effective rank
measures the number of significant dimensions in the representation;
the participation ratio is a complementary measure of dimensionality;
top-1 variance fraction measures how much of the total variance is
captured by the single leading principal component.

\begin{table}[h]
\centering
\begin{tabular}{lrrrr}
\toprule
\textbf{Layer} & \textbf{bp\_ce} & \textbf{bp\_kd} & \textbf{epc\_orig} & \textbf{epc\_deep} \\
\midrule
\multicolumn{5}{l}{\textit{Entropy Effective Rank}} \\
\midrule
embedding & 374.65 & 365.65 & 342.79 & 340.98 \\
block\_0 & 51.78 & 7.21 & 8.18 & 7.84 \\
block\_1 & 30.60 & 5.71 & 3.70 & 3.62 \\
block\_2 & 21.18 & 5.15 & 2.80 & 2.64 \\
block\_3 & 16.63 & 4.88 & 2.58 & 2.40 \\
block\_4 & 14.06 & 4.74 & 2.49 & 2.32 \\
block\_5 & 12.33 & 4.66 & 3.38 & 3.18 \\
final & 13.90 & 4.98 & 3.87 & 3.78 \\
\midrule
\multicolumn{5}{l}{\textit{Participation Ratio}} \\
\midrule
embedding & 186.77 & 169.86 & 129.88 & 126.88 \\
block\_0 & 13.70 & 3.43 & 4.15 & 3.99 \\
block\_1 & 9.83 & 3.21 & 2.38 & 2.32 \\
block\_2 & 7.77 & 3.12 & 1.91 & 1.80 \\
block\_3 & 6.65 & 3.05 & 1.80 & 1.67 \\
block\_4 & 5.95 & 3.00 & 1.77 & 1.63 \\
block\_5 & 5.42 & 2.98 & 2.28 & 2.15 \\
final & 6.14 & 3.13 & 2.62 & 2.59 \\
\midrule
\multicolumn{5}{l}{\textit{Top-1 Variance Fraction}} \\
\midrule
embedding & 0.0354 & 0.0418 & 0.0595 & 0.0610 \\
block\_0 & 0.1993 & 0.4480 & 0.4121 & 0.4230 \\
block\_1 & 0.2515 & 0.4710 & 0.6106 & 0.6211 \\
block\_2 & 0.2964 & 0.4773 & 0.7038 & 0.7294 \\
block\_3 & 0.3285 & 0.4876 & 0.7296 & 0.7626 \\
block\_4 & 0.3528 & 0.4981 & 0.7373 & 0.7734 \\
block\_5 & 0.3746 & 0.5033 & 0.6250 & 0.6475 \\
final & 0.3406 & 0.4899 & 0.5663 & 0.5685 \\
\midrule
\multicolumn{5}{l}{\textit{Dead Neuron Fraction}} \\
\midrule
all layers & 0.0000 & 0.0000 & 0.0000 & 0.0000 \\
final & 0.0013 & 0.0013 & 0.0013 & 0.0013 \\
\bottomrule
\end{tabular}
\caption{A1 representational collapse metrics (mean over 5 seeds).}
\label{tab:a1-collapse}
\end{table}

\textbf{Interpretation:}
EPC arms show dramatically more representational collapse.
At the final layer:
\begin{itemize}[nosep]
  \item Effective rank: bp\_ce 13.9, bp\_kd 5.0, epc\_original 3.9,
    epc\_deep 3.8.  EPC representations live in a $\sim$3.8-dimensional
    subspace.
  \item Top-1 variance: epc\_original 0.566, epc\_deep 0.569, vs
    bp\_ce 0.341.  A single component explains $\sim$57\% of EPC
    variance vs $\sim$34\% for CE.
  \item The collapse is not neuron death (dead fraction is identical
    at 0.13\% across all arms).  It is a low-dimensional
    \textit{representational} bottleneck.
  \item bp\_kd also shows significant compression (eff.\ rank 5.0),
    but not as severe as EPC.
  \item epc\_deep is slightly more collapsed than epc\_original,
    suggesting that increasing inference depth intensifies the
    compression rather than mitigating it.
\end{itemize}

%% ---- A2 ----
\subsection{A2: Factor Probe Accuracy}

Linear classifiers are trained on each layer's hidden states to
predict four syntactic factors.  Test accuracy on held-out sentences
measures whether the factor is linearly decodable; held-out accuracy
on unseen factor combinations measures compositional generalization.

\begin{table}[h]
\centering
\begin{tabular}{lrrrrrrrr}
\toprule
& \multicolumn{2}{c}{\textbf{subj\_num}} & \multicolumn{2}{c}{\textbf{obj\_num}} & \multicolumn{2}{c}{\textbf{tense}} & \multicolumn{2}{c}{\textbf{negation}} \\
\cmidrule(lr){2-3}\cmidrule(lr){4-5}\cmidrule(lr){6-7}\cmidrule(lr){8-9}
\textbf{Arm} & test & heldout & test & heldout & test & heldout & test & heldout \\
\midrule
\multicolumn{9}{l}{\textit{Final layer}} \\
\midrule
bp\_ce & 0.957 & 0.903 & 0.696 & 0.257 & 0.913 & 0.856 & 1.000 & 1.000 \\
bp\_kd & 0.838 & 0.732 & 0.688 & 0.120 & 0.824 & 0.763 & 0.987 & 0.978 \\
epc\_original & 0.750 & 0.658 & 0.671 & 0.139 & 0.771 & 0.769 & 0.983 & 0.984 \\
epc\_deep & 0.688 & 0.610 & 0.678 & 0.072 & 0.752 & 0.752 & 0.977 & 0.976 \\
\midrule
\multicolumn{9}{l}{\textit{Block 3}} \\
\midrule
bp\_ce & 0.998 & 0.995 & 0.769 & 0.457 & 0.991 & 0.973 & 1.000 & 1.000 \\
bp\_kd & 0.988 & 0.961 & 0.696 & 0.238 & 0.935 & 0.885 & 1.000 & 0.999 \\
epc\_original & 0.500 & 0.503 & 0.511 & 0.482 & 0.537 & 0.554 & 0.876 & 0.860 \\
epc\_deep & 0.506 & 0.516 & 0.565 & 0.340 & 0.637 & 0.640 & 0.880 & 0.874 \\
\bottomrule
\end{tabular}
\caption{A2 factor probe accuracy (mean over 5 seeds).  0.5 = chance.}
\label{tab:a2-factors}
\end{table}

\textbf{Interpretation:}
At block 3, EPC arms are at chance (0.50) for subj\_num, while bp\_ce
retains 0.998.  By the final layer, EPC partially recovers (0.69--0.75)
but still trails bp\_ce (0.957).  The same pattern holds for tense and
obj\_num.  Only negation is equally well-retained across all arms.
epc\_deep is slightly worse than epc\_original, again suggesting depth
intensifies information loss rather than helping.

%% ---- A3 ----
\subsection{A3: Causal Selectivity}

Causal selectivity measures the degree to which intervening on a
syntactic factor causally changes the representation, as measured by
CKA discrepancy.  Higher means more selective encoding of the factor.

\begin{table}[h]
\centering
\begin{tabular}{lrrrr}
\toprule
\textbf{Layer} & \textbf{bp\_ce} & \textbf{bp\_kd} & \textbf{epc\_orig} & \textbf{epc\_deep} \\
\midrule
\multicolumn{5}{l}{\textit{subj\_num selectivity}} \\
\midrule
embedding & 0.645 & 0.648 & 0.645 & 0.645 \\
block\_0 & 0.793 & 0.763 & 0.757 & 0.765 \\
block\_1 & 0.783 & 0.732 & 0.558 & 0.573 \\
block\_2 & 0.757 & 0.708 & 0.424 & 0.407 \\
block\_3 & 0.719 & 0.670 & 0.399 & 0.375 \\
block\_4 & 0.672 & 0.633 & 0.408 & 0.400 \\
block\_5 & 0.620 & 0.606 & 0.507 & 0.401 \\
final & 0.607 & 0.514 & 0.438 & 0.379 \\
\midrule
\multicolumn{5}{l}{\textit{obj\_num selectivity}} \\
\midrule
embedding & 0.670 & 0.666 & 0.663 & 0.663 \\
block\_0 & 0.578 & 0.507 & 0.531 & 0.530 \\
block\_1 & 0.557 & 0.491 & 0.348 & 0.352 \\
block\_2 & 0.518 & 0.468 & 0.283 & 0.265 \\
block\_3 & 0.496 & 0.453 & 0.343 & 0.350 \\
block\_4 & 0.461 & 0.404 & 0.358 & 0.344 \\
block\_5 & 0.425 & 0.357 & 0.301 & 0.254 \\
final & 0.425 & 0.302 & 0.232 & 0.275 \\
\midrule
\multicolumn{5}{l}{\textit{tense selectivity}} \\
\midrule
embedding & 0.719 & 0.726 & 0.730 & 0.730 \\
block\_0 & 0.834 & 0.809 & 0.784 & 0.780 \\
block\_1 & 0.813 & 0.766 & 0.618 & 0.610 \\
block\_2 & 0.789 & 0.731 & 0.552 & 0.510 \\
block\_3 & 0.749 & 0.704 & 0.557 & 0.509 \\
block\_4 & 0.707 & 0.657 & 0.490 & 0.448 \\
block\_5 & 0.674 & 0.627 & 0.566 & 0.511 \\
final & 0.661 & 0.592 & 0.546 & 0.481 \\
\midrule
\multicolumn{5}{l}{\textit{negation selectivity}} \\
\midrule
embedding & 0.685 & 0.680 & 0.676 & 0.676 \\
block\_0 & 0.891 & 0.868 & 0.863 & 0.874 \\
block\_1 & 0.899 & 0.863 & 0.764 & 0.796 \\
block\_2 & 0.908 & 0.848 & 0.662 & 0.712 \\
block\_3 & 0.909 & 0.815 & 0.581 & 0.592 \\
block\_4 & 0.904 & 0.767 & 0.569 & 0.566 \\
block\_5 & 0.887 & 0.713 & 0.749 & 0.754 \\
final & 0.870 & 0.720 & 0.719 & 0.718 \\
\bottomrule
\end{tabular}
\caption{A3 causal selectivity scores (mean over 5 seeds).}
\label{tab:a3-selectivity}
\end{table}

\textbf{Interpretation:}
All arms start at similar selectivity at the embedding layer (0.64--0.73
across factors), confirming the input representation is identical.
The divergence begins at block 1 and widens through the depth of the
network:

\begin{itemize}[nosep]
  \item \textbf{subj\_num at block 3:} bp\_ce 0.719, bp\_kd 0.670,
    epc\_original 0.399, epc\_deep 0.375.  EPC loses more than half
    the selectivity signal by mid-network.
  \item \textbf{tense at block 3:} bp\_ce 0.749, epc\_original 0.557.
    Same trend.
  \item \textbf{negation at block 3:} bp\_ce 0.909, epc\_original
    0.581.  Even the most robust factor degrades sharply.
  \item \textbf{epc\_deep is consistently $\leq$ epc\_original} on
    selectivity, confirming that deeper inference intensifies
    information loss.
\end{itemize}

%% ================================================================
\section{Synthesis and Interpretation}
%% ================================================================

\subsection{Convergent Evidence}

Three independent lines of evidence converge on the same conclusion:

\begin{enumerate}
  \item \textbf{R8 Performance:} EPC does not beat baselines on raw
    validation loss at matched updates (6.75--6.78 vs 6.26 for CE)
    and is dominated at matched wall-clock (327 vs 863--889 perplexity).

  \item \textbf{R9-A1 Collapse:} EPC representations are severely
    compressed (final effective rank 3.8--3.9 vs 13.9 for CE).  A
    single component captures 57\% of variance.

  \item \textbf{R9-A2/A3 Factors:} EPC loses syntactic factor
    information faster than baselines.  At block 3, factor probes are
    at chance for subj\_num (0.50) and causal selectivity drops to
    0.38--0.40 vs 0.72 for CE.
\end{enumerate}

\subsection{Mechanistic Read}

The credit assignment diagnostic (Table~\ref{tab:credit-assignment})
shows that the ePC gradient signal is effectively zero for blocks 0--4
and only non-zero at the final block (h.5), with a norm ratio of
$\sim$$3.3 \times 10^{-5}$ relative to ordinary KD.  This means the
energy--monotonicity constraint is structurally enforced but its
credit signal is negligible for the lower layers.

The most likely mechanism is that ePC acts as an excessively strong
attractor/regularizer: by enforcing energy monotonicity through
iterative inference, it drives the network toward a low-energy
fixed point that compresses the representation space.  The constraint
is satisfied (hence \texttt{activity\_energy\_monotone: true}), but at
the cost of discarding task-relevant degrees of freedom.

The fact that epc\_deep ($T{=}12$) is consistently equal to or
slightly worse than epc\_original on every metric --- performance,
collapse, factor retention, and selectivity --- confirms that
increasing inference depth does not reveal a hidden advantage.  It
intensifies the compression.

\subsection{What Is Validated vs What Is Not}

\textbf{Validated:}
\begin{itemize}[nosep]
  \item The energy--monotonicity property is real and robustly
    enforceable (10/10 EPC records across 5 seeds).
  \item The probe infrastructure is trustworthy (A0 calibration passes
    at machine precision).
\end{itemize}

\textbf{Not supported:}
\begin{itemize}[nosep]
  \item EPC does not improve predictive performance over plain
    baselines.
  \item EPC does not produce cleaner or more disentangled
    representations; it produces more collapsed ones.
  \item EPC does not preserve causal factor information better; it
    loses it faster.
  \item Increasing inference depth does not help.
\end{itemize}

\textbf{Caveats:}
\begin{itemize}[nosep]
  \item These results are for a 6-layer GPT-2 student on a synthetic
    grammar corpus.  Whether the pattern holds at larger scale or on
    natural language is not tested.
  \item Stage A probes measure linear decodability and causal
    selectivity, not all possible representational properties.
  \item The ePC credit signal being concentrated in the final block
    suggests the current implementation may not be distributing credit
    as intended; a modified formulation could behave differently.
\end{itemize}

\subsection{Recommendation}

The current ePC formulation does not warrant advancement toward an
adaptive residual layer.  The structural property is valid but the
practical effect is counterproductive: it compresses representations
and discards factor information without a compensating performance or
generalization benefit.

If the energy--monotonicity idea is to be pursued further, the next
step should be diagnosing why the ePC credit signal is negligible for
all but the final block, and exploring formulations that distribute
credit more broadly before investing in additional scale or depth.

%% ================================================================
\section{Provenance}
%% ================================================================

\begin{itemize}[nosep]
  \item \textbf{R8 experiment:}
    \texttt{projects/relaleap/experiments/20260720T193211Z-epc-r8-broad-outcome/}
  \item \textbf{R8 artifacts:}
    \texttt{.../artifacts/r8/results/r8/train/training\_summary.json}
    (25 records, schema \texttt{relaleap.epc\_r8\_train.v1})
  \item \textbf{R9 Stage A experiment:}
    \texttt{projects/relaleap/experiments/20260721T222400Z-r9-stage-a/}
  \item \textbf{R9 artifacts:}
    \texttt{.../artifacts/seed*.json} (20 per-seed files) +
    \texttt{summary.json}
  \item \textbf{Source commit:} \texttt{ecf2f79}
  \item \textbf{Pod:} RunPod \texttt{0i26hl3ls399x4} (relaleap-r8-r9),
    A100-SXM4-80GB, \$1.49/hr
  \item \textbf{Seeds:} 1729, 3253, 6421, 8191, 10103
  \item \textbf{Teacher:} GPT-2 small (val\_loss 4.124, val\_ppl 61.8)
  \item \textbf{Student:} 6-layer GPT-2
  \item \textbf{Corpus:} synthetic subject--verb--object grammar
  \item \textbf{Report date:} 2026-07-22
\end{itemize}

\end{document}
